\documentclass[11pt,a4paper]{article}

\usepackage[utf8]{inputenc}
\usepackage[T1]{fontenc}
\usepackage{amsmath,amssymb,amsthm}
\usepackage{mathtools}
\usepackage{listings}
\usepackage{xcolor}
\usepackage{hyperref}
\usepackage{geometry}
\usepackage{booktabs}
\usepackage{enumitem}

\geometry{margin=2.5cm}

% Lean 4 listing style
\lstdefinelanguage{Lean4}{
  keywords={theorem,def,lemma,private,by,where,let,if,then,else,match,with,
            have,show,exact,intro,obtain,rw,simp,omega,subst,rcases,
            calc,refine,cases,apply,exfalso,push_neg,unfold,conv_lhs,
            abbrev,inductive,deriving,structure,import,open,section,
            namespace,end,instance,class,extends,fun,return,do},
  keywordstyle=\color{blue}\bfseries,
  comment=[l]{--},
  morecomment=[s]{/-}{-/},
  commentstyle=\color{gray}\itshape,
  stringstyle=\color{red!70!black},
  morestring=[b]",
  sensitive=true,
  literate=
    {forall}{{$\forall$}}1
    {exists}{{$\exists$}}1
    {=>}{{$\Rightarrow$}}2
    {->}{{$\to$}}2
    {:=}{{$\coloneqq$}}2
    {<=}{{$\le$}}2
    {>=}{{$\ge$}}2
    {/\\}{{$\land$}}2
    {\\/}{{$\lor$}}2
    {!=}{{$\ne$}}2
    {|-}{{$\vdash$}}2
    {Fin}{{$\mathsf{Fin}$}}3
    {Nat}{{$\mathbb{N}$}}3
    {Vertex}{{$\mathsf{Vertex}$}}6
    {Dir}{{$\mathsf{Dir}$}}3,
  basicstyle=\ttfamily\small,
  breaklines=true,
  frame=single,
  framerule=0.4pt,
  backgroundcolor=\color{gray!5},
  numbers=left,
  numberstyle=\tiny\color{gray},
  xleftmargin=1.5em,
  framexleftmargin=1.5em,
}

\lstset{language=Lean4}

\newtheorem{theorem}{Theorem}[section]
\newtheorem{lemma}[theorem]{Lemma}
\newtheorem{definition}[theorem]{Definition}
\newtheorem{remark}[theorem]{Remark}

\title{A Lean~4 Formalization of Knuth's ``Claude's Cycles''}
\author{Formalized in Lean~4 with Mathlib}
\date{March 2026}

\begin{document}

\maketitle

\begin{abstract}
We present a complete Lean~4 formalization (3{,}389 lines, zero \texttt{sorry})
of the main theorem from Knuth's ``Claude's Cycles'' (2026):
for every odd integer $m \ge 3$, the Cayley digraph on $m^3$ vertices
with arcs $ijk \to (i{+}1)jk$, $i(j{+}1)k$, $ij(k{+}1)$ (mod~$m$)
admits a decomposition of all $3m^3$ arcs into three directed Hamiltonian cycles.
The proof is carried out entirely within Lean~4's dependent type theory,
relying on Mathlib for modular arithmetic and finite-type reasoning.
\end{abstract}

\tableofcontents
\newpage

%% ============================================================
\section{Introduction}
\label{sec:intro}

Knuth's paper~\cite{knuth2026} studies a Cayley digraph $\mathcal{G}_m$ whose vertex set is
$\{0, 1, \ldots, m{-}1\}^3$ and whose arcs consist of incrementing any single
coordinate modulo~$m$.  Each vertex has out-degree~3, so the digraph has $3m^3$ arcs.
The main result states that for odd $m \ge 3$, these arcs partition into three directed
Hamiltonian cycles---cycles that visit all $m^3$ vertices exactly once.

The three cycles are defined by \emph{direction functions} $\mathsf{dir}_0$,
$\mathsf{dir}_1$, $\mathsf{dir}_2$ that assign to each vertex which coordinate
to increment next.  The key properties to verify are:
\begin{enumerate}[label=(\roman*)]
  \item At every vertex, the three directions form a permutation of $\{+i, +j, +k\}$.
  \item Each direction function generates a single cycle of length~$m^3$.
\end{enumerate}

Property~(i) ensures the arcs are partitioned; property~(ii) ensures each
piece is a Hamiltonian cycle.

%% ============================================================
\section{Formalization Framework}
\label{sec:framework}

\subsection{Types and Definitions}

Vertices are represented as triples of elements of $\mathsf{Fin}\;m$:

\begin{lstlisting}[caption={Basic types}]
abbrev Vertex (m : Nat) := Fin m * Fin m * Fin m

inductive Dir where
  | bumpi : Dir
  | bumpj : Dir
  | bumpk : Dir
\end{lstlisting}

The bump function increments a single coordinate:

\begin{lstlisting}[caption={Bump and orbit}]
def bump (m : Nat) [NeZero m] (v : Vertex m) (d : Dir) : Vertex m :=
  match d with
  | Dir.bumpi => (v.1 + 1, v.2.1, v.2.2)
  | Dir.bumpj => (v.1, v.2.1 + 1, v.2.2)
  | Dir.bumpk => (v.1, v.2.1, v.2.2 + 1)

def successor (m : Nat) [NeZero m]
    (dirFn : Vertex m -> Dir) (v : Vertex m) : Vertex m :=
  bump m v (dirFn v)

def orbit (m : Nat) [NeZero m]
    (dirFn : Vertex m -> Dir) (start : Vertex m) : (n : Nat) -> Vertex m
  | 0 => start
  | n + 1 => successor m dirFn (orbit m dirFn start n)
\end{lstlisting}

\subsection{The Fiber Index}

A central invariant is the \emph{fiber index} $s = (i + j + k) \bmod m$.

\begin{lstlisting}[caption={Fiber index}]
def fiberIndex (m : Nat) [NeZero m] (v : Vertex m) : Fin m :=
  v.1 + v.2.1 + v.2.2
\end{lstlisting}

\begin{lemma}[Fiber advancement]
For any vertex $v$ and direction $d$,
$\mathsf{fiberIndex}(\mathsf{bump}\;v\;d) = \mathsf{fiberIndex}(v) + 1$.
\end{lemma}

\begin{lstlisting}[caption={Fiber advances}]
theorem fiber_advances (m : Nat) [NeZero m] (v : Vertex m) (d : Dir) :
    fiberIndex m (bump m v d) = fiberIndex m v + 1 := by
  cases d <;> simp [fiberIndex, bump, add_assoc, add_comm, add_left_comm]
\end{lstlisting}

This means that within each round of $m$ steps, the fiber index cycles through all
$m$ residues. The step at which $s = 0$ marks the ``start'' of a new round.

%% ============================================================
\section{Direction Functions}
\label{sec:directions}

The three direction functions are decision tables based on the fiber index $s$
and coordinates $i$, $j$:

\begin{lstlisting}[caption={The three direction functions}]
def dir0 (m : Nat) [NeZero m] (v : Vertex m) : Dir :=
  let s := fiberIndex m v; let i := v.1; let j := v.2.1
  if s.val = 0 then (if j.val = m - 1 then Dir.bumpi else Dir.bumpk)
  else if s.val = m - 1 then (if i.val > 0 then Dir.bumpj else Dir.bumpk)
  else (if i.val = m - 1 then Dir.bumpk else Dir.bumpj)

def dir1 (m : Nat) [NeZero m] (v : Vertex m) : Dir :=
  let s := fiberIndex m v; let i := v.1
  if s.val = 0 then Dir.bumpj
  else if s.val = m - 1 then
    (if i.val > 0 then Dir.bumpk else Dir.bumpj)
  else Dir.bumpi

def dir2 (m : Nat) [NeZero m] (v : Vertex m) : Dir :=
  let s := fiberIndex m v; let i := v.1; let j := v.2.1
  if s.val = 0 then (if j.val < m - 1 then Dir.bumpi else Dir.bumpk)
  else if s.val = m - 1 then Dir.bumpi
  else (if i.val < m - 1 then Dir.bumpk else Dir.bumpj)
\end{lstlisting}

\begin{theorem}[Direction permutation]
At every vertex, $\{\mathsf{dir}_0(v),\;\mathsf{dir}_1(v),\;\mathsf{dir}_2(v)\}
= \{+i, +j, +k\}$.
\end{theorem}

\begin{lstlisting}[caption={Permutation property}]
theorem directions_are_permutation (m : Nat) [NeZero m] (hm : m >= 3)
    (v : Vertex m) :
    ({dir0 m v, dir1 m v, dir2 m v} : Finset Dir)
      = {Dir.bumpi, Dir.bumpj, Dir.bumpk} := by
  simp only [dir0, dir1, dir2]
  have hj_bound := v.2.1.isLt; have hi_bound := v.1.isLt
  split_ifs <;> first
    | (ext x; fin_cases x <;> simp_all <;> done) | omega
\end{lstlisting}

%% ============================================================
\section{Cycle 0: Block Decomposition}
\label{sec:cycle0}

\subsection{Strategy}

The orbit under $\mathsf{dir}_0$ is analyzed by decomposing the $m^3$ vertices into
$m$ \emph{blocks} $B_0, B_1, \ldots, B_{m-1}$, where $B_i$ contains all vertices
with first coordinate equal to~$i$.  Each block has $m^2$ vertices.

The proof proceeds in three stages:
\begin{enumerate}
  \item \textbf{Intra-block completeness}: Within each block, the orbit visits
    all $m^2$ vertices via a closed-form formula.
  \item \textbf{Block transitions}: At the exit of block~$i$, the direction is
    $\mathsf{bumpi}$, which transitions the orbit to block~$i{+}1$.
  \item \textbf{Chaining}: After $m$ blocks, the orbit returns to the start,
    completing a cycle of length $m \cdot m^2 = m^3$.
\end{enumerate}

\subsection{Closed-Form Formulas}

For the $i = 0$ block, the expected vertex at orbit step $t \cdot m + r$
(where $0 \le t, r < m$) is:

\begin{lstlisting}[caption={Expected formula for the $i=0$ block}]
def i0Expected (m : Nat) [NeZero m] (t r : Fin m) : Vertex m :=
  let j_base : Fin m := { val := ((m - 2) * t.val) % m, ... }
  let k_base : Fin m := { val := (2 * t.val) % m, ... }
  if r.val = 0 then
    ({ val := 0, ... }, j_base, k_base)
  else
    ({ val := 0, ... }, j_base + { val := r.val - 1, ... }, k_base + 1)
\end{lstlisting}

The orbit-matching theorem is proved by strong induction on the step count,
with a case split on whether $r = 0$ (start of round), $0 < r < m-1$ (mid-round),
or $r = m-1$ (end of round):

\begin{lstlisting}[caption={Orbit matches the formula}]
theorem orbit_i0_explicit (m : Nat) [NeZero m]
    (hm_odd : m % 2 = 1) (hm_ge : m >= 3) (t r : Fin m)
    (hn : n = t.val * m + r.val) :
    orbit m (dir0 m) (checkpoint0 m { val := 0, ... }) n
      = i0Expected m t r
\end{lstlisting}

Intermediate blocks ($0 < i < m{-}1$) and the last block ($i = m{-}1$)
follow the same pattern with minor variations in the coordinate formulas.

\subsection{Injectivity via Coprimality}

\begin{lemma}[Coprimality cancellation]
\label{lem:coprime}
If $m$ is odd and $2t_1 \equiv 2t_2 \pmod{m}$ with $0 \le t_1, t_2 < m$,
then $t_1 = t_2$.
\end{lemma}

\begin{lstlisting}[caption={Coprimality cancellation}]
private theorem t_eq_of_two_t_mod (m : Nat) (hm_odd : m % 2 = 1) (hm_ge : m >= 3)
    (t1 t2 : Fin m) (h : (2 * t1.val) % m = (2 * t2.val) % m) : t1 = t2 := by
  by_contra hne
  have hne_val : t1.val != t2.val := fun heq => hne (Fin.ext heq)
  rcases Nat.lt_or_gt_of_ne hne_val with hlt | hgt
  -- t1 < t2: m | 2*(t2-t1), bounded by 2m, so 2*(t2-t1) = m => m even
  . have hdvd : m | (2 * t2.val - 2 * t1.val) :=
      (Nat.modEq_iff_dvd' (by omega)).mp h
    obtain <<c, hc>> := hdvd
    have : c = 1 := by nlinarith [t1.isLt, t2.isLt]
    subst this; simp only [Nat.mul_one] at hc; omega
  -- t1 > t2: symmetric
  . ...
\end{lstlisting}

The proof exploits that $m \mid 2(t_2 - t_1)$ with $0 < 2(t_2 - t_1) < 2m$
forces $2(t_2 - t_1) = m$, contradicting $m$ odd.

This lemma is used to recover $t$ from the coordinates:
the $j$-base and $k$-base coordinates involve $((m{-}2) \cdot t) \bmod m$
and $(2t) \bmod m$, and coprimality ensures distinct $t$ values produce distinct sums.

%% ============================================================
\section{Cycles 1 and 2: Super-Round Architecture}
\label{sec:cycles12}

\subsection{Uniform Proof Pattern}

Cycles~1 and~2 follow a common architecture.  The $m^3$ orbit steps are grouped into
$m$ \emph{super-rounds}, each of $m^2$ steps, indexed by a parameter $u \in \{0, \ldots, m{-}1\}$.
Within super-round~$u$, steps are indexed as $t \cdot m + r$.

For each cycle, we define:
\begin{itemize}
  \item A \textbf{block entry} function $\mathsf{BlockEntry}(u)$: the vertex where
    super-round~$u$ begins.
  \item An \textbf{Expected} function $\mathsf{Expected}(u, t, r)$: the closed-form vertex
    at step $u \cdot m^2 + t \cdot m + r$.
\end{itemize}

Hamiltonicity is then proved via four components:

\begin{center}
\begin{tabular}{ll}
\toprule
\textbf{Component} & \textbf{Statement} \\
\midrule
\texttt{orbit\_explicit} & $\mathsf{orbit}^{t \cdot m + r}(\mathsf{BlockEntry}(u)) = \mathsf{Expected}(u, t, r)$ \\
\texttt{Expected\_injective} & $(u_1, t_1, r_1) \ne (u_2, t_2, r_2) \Rightarrow \mathsf{Expected}(\cdots) \ne \mathsf{Expected}(\cdots)$ \\
\texttt{Expected\_surjective} & Follows from injectivity on a finite type \\
\texttt{orbit\_returns} & $\mathsf{orbit}^{m^3}(\mathsf{BlockEntry}(0)) = \mathsf{BlockEntry}(0)$ \\
\bottomrule
\end{tabular}
\end{center}

\subsection{Cycle 1: Details}

The expected formula for cycle~1 has two branches:

\begin{lstlisting}[caption={Cycle 1 expected formula}]
def dir1Expected (m : Nat) [NeZero m] (u t r : Fin m) : Vertex m :=
  let i_base := { val := ((m - 2) * t.val) % m, ... }
  let j_base := { val := (u.val + t.val) % m, ... }
  let k_val  := { val := (t.val + m - u.val) % m, ... }
  if r.val = 0 then (i_base, j_base, k_val)
  else (i_base + { val := r.val - 1, ... }, j_base + 1, k_val)
\end{lstlisting}

Injectivity recovery strategy:
\begin{enumerate}
  \item Recover $r$ from the fiber index $s = (i + j + k) \bmod m$.
  \item Recover $t$ from $j + k = (2t + m) \bmod m$ using Lemma~\ref{lem:coprime}.
  \item Recover $u$ from $j - t = u \pmod{m}$.
\end{enumerate}

\subsection{Cycle 2: The Most Complex Proof}

The expected formula for cycle~2 has \emph{five} branches, depending on whether the
super-round is ``non-special'' ($j_u = ((m{-}2) \cdot u) \bmod m < m{-}1$)
or ``special'' ($j_u = m{-}1$):

\begin{lstlisting}[caption={Cycle 2 expected formula (simplified)}]
def dir2Expected (m : Nat) [NeZero m] (u t r : Fin m) : Vertex m :=
  let j_u := ((m - 2) * u.val) % m
  if j_u < m - 1 then
    -- Non-special: 3 sub-branches (r=0, r>=1 with t<m-1, r>=1 with t=m-1)
    if r.val = 0 then (2t mod m, j_u, (2u + (m-2)t) mod m)
    else if t.val < m - 1 then (2t+1 mod m, j_u, (2u + (m-2)t + r - 1) mod m)
    else (m-1, (j_u + r - 1) mod m, (2u + (m-2)(m-1)) mod m)
  else
    -- Special: 2 sub-branches (t<m-1 or r<=1, otherwise)
    if t.val < m - 1 \/ r.val <= 1 then (t, m-1, (2u + m - t + r) mod m)
    else (m-1, (r-2) mod m, (2u + 2) mod m)
\end{lstlisting}

\subsubsection{The \texttt{u\_eq\_of\_m2\_u\_mod} Helper}

A key helper reduces cancellation of $(m{-}2)$ to cancellation of~$2$:

\begin{lstlisting}[caption={Reducing $(m{-}2)$-cancellation to $2$-cancellation}]
private theorem u_eq_of_m2_u_mod (m : Nat) (hm_odd : m % 2 = 1) (hm_ge : m >= 3)
    (u1 u2 : Fin m) (h : ((m - 2) * u1.val) % m = ((m - 2) * u2.val) % m)
    : u1 = u2 := by
  have h2u : (2 * u1.val) % m = (2 * u2.val) % m := by
    -- Key identity: (m-2)*u + 2*u = m*u, so (m-2)*u = -2*u (mod m)
    have hfact1 : (m - 2) * u1.val + 2 * u1.val = m * u1.val := by
      zify [show m >= 2 from by omega]; ring
    ...
  exact t_eq_of_two_t_mod m hm_odd hm_ge u1 u2 h2u
\end{lstlisting}

\subsubsection{Injectivity: 25-Case Analysis}

The injectivity proof for \texttt{dir2Expected} involves 25 sub-cases from the
$4 \times$ sub-case structure:

\begin{enumerate}
  \item \textbf{Case A} (both non-special): 9 sub-cases from $(r_1{=}0$ vs $r_1{\ge}1) \times (t_1{<}m{-}1$ vs $t_1{=}m{-}1) \times (\text{same for } u_2, t_2)$.
  \item \textbf{Case B} (non-special vs special): Contradiction from the $j$-component.
  \item \textbf{Case C} (special vs non-special): Symmetric to Case~B.
  \item \textbf{Case D} (both special): $u_1 = u_2$ from $j_u$, then $t$ from bounds.
\end{enumerate}

The proof structure in Lean:

\begin{lstlisting}[caption={Injectivity proof structure}]
private theorem dir2Expected_injective (m : Nat) [NeZero m]
    (hm_odd : m % 2 = 1) (hm_ge : m >= 3) :
    Function.Injective (fun p : Fin m * Fin m * Fin m =>
      dir2Expected m p.1 p.2.1 p.2.2) := by
  intro <<u1, t1, r1>> <<u2, t2, r2>> heq
  -- Step 1: r1 = r2 from fiber index
  have hr : r1 = r2 := ... -- via dir2E_fiber
  subst hr
  -- Step 2: Unfold and case-split on special/non-special
  simp only [dir2Expected] at heq
  by_cases hns1 : ((m - 2) * u1.val) % m < m - 1 <;>
  by_cases hns2 : ((m - 2) * u2.val) % m < m - 1
  . -- Case A: Both non-special (9 sub-cases)
    ...
  . -- Case B: u1 non-special, u2 special -> contradiction
    exfalso; ...
  . -- Case C: u1 special, u2 non-special -> contradiction
    exfalso; ...
  . -- Case D: Both special
    ...
\end{lstlisting}

\subsubsection{Omega Limitations and Workarounds}

A recurring technical challenge: Lean's \texttt{omega} tactic handles Presburger arithmetic
but \emph{cannot} reason about modular arithmetic with variable modulus.
For example, \texttt{omega} cannot prove:
\[
  m \le 2t + 1 < 2m \implies (2t + 1) \bmod m = 2t + 1 - m
\]
The workaround decomposes the modular equation explicitly:

\begin{lstlisting}[caption={Working around omega's limitations}]
have : (2 * t1.val + 1) % m = 2 * t1.val + 1 - m := by
  conv_lhs => rw [show 2 * t1.val + 1
    = (2 * t1.val + 1 - m) + m from by omega]
  rw [Nat.add_mod_right, Nat.mod_eq_of_lt (by omega)]
\end{lstlisting}

Another subtlety: \texttt{Nat.ModEq m a b} is definitionally $a \bmod m = b \bmod m$
but displayed as $a \equiv b \pmod{m}$, and \texttt{rw} cannot see through the notation:

\begin{lstlisting}[caption={Unfolding Nat.ModEq}]
have hju := Nat.ModEq.add_right_cancel' (r1.val - 1) hj
unfold Nat.ModEq at hju  -- Must unfold before rewriting
rw [Nat.mod_eq_of_lt (Nat.mod_lt _ hm_pos),
    Nat.mod_eq_of_lt (Nat.mod_lt _ hm_pos)] at hju
\end{lstlisting}

%% ============================================================
\section{Assembly of the Main Theorem}
\label{sec:assembly}

\subsection{Individual Hamiltonicity Theorems}

Each cycle's Hamiltonicity theorem has the same form:

\begin{lstlisting}[caption={Cycle Hamiltonicity (pattern)}]
theorem cycleN_hamiltonian (m : Nat) [NeZero m]
    (hm_odd : m % 2 = 1) (hm_ge : m >= 3) :
    exists start : Vertex m,
      (orbit m (dirN m) start (m ^ 3) = start) /\
      (forall v : Vertex m, exists n : Fin (m ^ 3),
        orbit m (dirN m) start n.val = v)
\end{lstlisting}

The proof instantiates the start vertex, appeals to orbit return for cyclicity,
and appeals to surjectivity (from injectivity on a finite type) for vertex coverage:

\begin{lstlisting}[caption={Surjectivity from injectivity on finite types}]
private theorem dir2Expected_surjective ... :
    Function.Surjective (fun p : Fin m * Fin m * Fin m =>
      dir2Expected m p.1 p.2.1 p.2.2) :=
  Finite.surjective_of_injective (dir2Expected_injective m hm_odd hm_ge)
\end{lstlisting}

\subsection{The Decomposition Theorem}

The final assembly is a four-component conjunction:

\begin{lstlisting}[caption={The main decomposition theorem}]
theorem claude_decomposition (m : Nat) [NeZero m]
    (hm_odd : m % 2 = 1) (hm_ge : m >= 3) :
    (exists start0, orbit m (dir0 m) start0 (m ^ 3) = start0 /\
      forall v, exists n : Fin (m ^ 3),
        orbit m (dir0 m) start0 n.val = v) /\
    (exists start1, orbit m (dir1 m) start1 (m ^ 3) = start1 /\
      forall v, exists n : Fin (m ^ 3),
        orbit m (dir1 m) start1 n.val = v) /\
    (exists start2, orbit m (dir2 m) start2 (m ^ 3) = start2 /\
      forall v, exists n : Fin (m ^ 3),
        orbit m (dir2 m) start2 n.val = v) /\
    (forall v : Vertex m,
      ({dir0 m v, dir1 m v, dir2 m v} : Finset Dir)
        = {Dir.bumpi, Dir.bumpj, Dir.bumpk}) := by
  exact <<
    cycle0_hamiltonian m hm_odd hm_ge,
    cycle1_hamiltonian m hm_odd hm_ge,
    cycle2_hamiltonian m hm_odd hm_ge,
    directions_are_permutation m hm_ge >>
\end{lstlisting}

%% ============================================================
\section{Proof Statistics and Structure}
\label{sec:stats}

\begin{table}[h]
\centering
\begin{tabular}{lrrr}
\toprule
\textbf{Component} & \textbf{Lines} & \textbf{Theorems} & \textbf{Max case depth} \\
\midrule
Basic definitions \& properties     & 130   & 12  & --- \\
Cycle 0 (block analysis)            & 1{,}595 & 45 & 3 \\
Cycle 0 Hamiltonicity assembly      & 55    & 2   & --- \\
Cycle 1 (super-round analysis)      & 580   & 25  & 9 \\
Cycle 2 (super-round analysis)      & 980   & 28  & 25 \\
Decomposition theorem               & 20    & 1   & --- \\
\midrule
\textbf{Total}                      & \textbf{3{,}389} & \textbf{$\sim$110} & \\
\bottomrule
\end{tabular}
\caption{Formalization statistics by section.}
\end{table}

\begin{table}[h]
\centering
\begin{tabular}{lr}
\toprule
\textbf{Metric} & \textbf{Value} \\
\midrule
Total lines of Lean code & 3{,}389 \\
\texttt{sorry} count & 0 \\
Mathlib imports & 5 \\
Approximate build time & $\sim$100\,s \\
Lean version & 4.29.0-rc2 \\
\bottomrule
\end{tabular}
\caption{Overall metrics.}
\end{table}

%% ============================================================
\section{Key Proof Techniques}
\label{sec:techniques}

\subsection{Coprimality Arguments}

Since $\gcd(2, m) = 1$ for odd~$m$, the map $t \mapsto 2t \bmod m$ is a bijection
on $\{0, \ldots, m{-}1\}$.  This is formalized as \texttt{t\_eq\_of\_two\_t\_mod}
(Section~\ref{sec:cycle0}).  The identity $(m{-}2) \cdot u + 2u = mu$ further yields
$(m{-}2) \cdot u \equiv -2u \pmod{m}$, giving coprimality for $(m{-}2)$ as well.

\subsection{Fiber Index for Parameter Recovery}

Each direction function $\mathsf{dir}_c$ is designed so that within a round of $m$ steps,
the fiber index $s$ cycles through all residues $0, 1, \ldots, m{-}1$.  This means the
intra-round step~$r$ can be recovered from any vertex via $r = s \bmod m$
(or a derived quantity).

\subsection{Injectivity $\Rightarrow$ Surjectivity on Finite Types}

The key insight that simplifies the proof: on a finite type, injectivity implies
surjectivity.  Mathlib provides \texttt{Finite.surjective\_of\_injective} for this.
Once we show $\mathsf{Expected}$ is injective, surjectivity (and hence coverage of all
$m^3$ vertices) is automatic.

\subsection{Dealing with Lean~4 Technicalities}

Several Lean-specific challenges arose:
\begin{itemize}
  \item \textbf{If-then-else and projections}: Tuple projections like \texttt{.2.1}
    do not reduce through unresolved if-then-else expressions.
    The fix: resolve all conditional branches with
    \texttt{simp only [if\_pos h, if\_neg h]} \emph{before} destructuring.
  \item \textbf{Dependent elimination}: \texttt{obtain $\langle$a, b, c$\rangle$ := heq} fails
    when \texttt{heq} contains \texttt{Decidable.rec} from unresolved conditionals.
  \item \textbf{Variable representation after \texttt{subst}}: After substitution,
    Lean may represent $t_1$ as \texttt{(u1, t1, r1).2.1} internally.
    The \texttt{omega} tactic then treats the hypothesis \texttt{t1.val} and the
    goal's \texttt{(u1, t1, r1).2.1.val} as distinct variables.
    Workaround: prove equalities with \texttt{Fin.ext} in a separate \texttt{have},
    then \texttt{subst}.
\end{itemize}

%% ============================================================
\section{Conclusion}

We have presented a complete, machine-checked proof that for odd $m \ge 3$,
the Cayley digraph $\mathcal{G}_m$ admits a decomposition of its $3m^3$ arcs into
three directed Hamiltonian cycles.  The formalization closely follows Knuth's
construction but required significant effort in managing modular arithmetic
within Lean~4's type theory, particularly the interplay between
$\mathsf{Fin}\;m$ arithmetic and natural number modular operations.

The most technically demanding component was the injectivity of the 5-branch
formula for cycle~2 (\texttt{dir2Expected\_injective}, 220~lines), which
required 25 sub-cases and multiple workarounds for limitations of the
\texttt{omega} tactic with variable-modulus arithmetic.

\begin{thebibliography}{9}
\bibitem{knuth2026}
  D.~E.~Knuth,
  \textit{Claude's Cycles},
  Manuscript, 28~February 2026 (revised 02~March 2026).
\end{thebibliography}

\end{document}
